% Sections 4.1 to 4.9, from lectures/L04/appendix/a_metastability_and_synchronization.md: one
% section per part, A.1 to A.9.

\section{The Problem: A Button Isn't a Testbench Signal}
\label{c4:a1}
In Chapter~\ref{c3:ch} you built edge detectors and reasoned about them against a clean clock, where
every input changed neatly between edges, exactly the way a simulator draws it. A push button, an
external reset line, or any signal from another clock domain changes whenever the outside world
decides to, completely independent of your FPGA's clock.

Every D flip-flop has two timing requirements around the active clock edge:
\begin{itemize}
  \item a \textbf{setup time}: how long \code{D} must be stable \emph{before} the edge.
  \item a \textbf{hold time}: how long \code{D} must stay stable \emph{after} it.
\end{itemize}
If an input changes inside that window the flip-flop is not guaranteed to capture a clean \code{0}
or \code{1}, and you can never promise it will land outside, because you do not control when an
external signal changes. The rest of this chapter solves that not by avoiding the problem but by
containing it.

\section{Metastability: What Actually Happens Inside the Flip-Flop}
\label{c4:a2}
A D flip-flop is built from cross-coupled gates that snap to \code{0} or \code{1} the instant the
clock samples the input. If \code{D} changes too close to the edge, the internal feedback loop can
be left balanced between the two states, neither fully high nor fully low, for longer than
expected. That is \textbf{metastability}.

Picture a ball balanced exactly on top of a hill. Given enough time it rolls down one side, but
exactly when it commits, and to which side, is not predictable from the moment it was placed there,
and while it is still balanced its position is not a valid digital value at all.

A metastable output behaves the same way:
\begin{itemize}
  \item it almost always resolves to a valid \code{0} or \code{1} on its own, given enough time.
  \item ``enough time'' is not free, and is not the same on every clock cycle.
  \item if the unresolved voltage reaches downstream logic before it settles, different gates can
    legitimately read it differently: one sees \code{0}, another \code{1}, at the same instant.
\end{itemize}
That disagreement is where real systems misbehave: a counter that increments in one part of the
circuit but not another, or a state machine in an impossible state. Everything in this chapter
exists to make sure metastability resolves \emph{before} a signal touches the rest of your design.

In time, the whole problem looks like this:

\bookfigure{setup_hold}{\code{D} changes inside the setup/hold window around a rising clock edge,
leaving \code{Q} metastable before it settles, late, to either level.}{c4:fig:setup_hold}

\noindent Read it left to right. The shaded band is the window from \secref{c4:a1}, \code{t_su}
before the edge and \code{t_h} after it. \code{D} changes inside it, which is exactly what an input
you do not control is free to do. \code{Q} is then neither \code{0} nor \code{1} for a while, and
when it does settle, \emph{when} it settles and \emph{which level} it settles to are both outside
your control. Both dashed continuations are legitimate outcomes of the same input.

Note what the diagram does not show: a wrong answer. Sampling a signal mid-transition is allowed to
give you either level, and that on its own is harmless, because the input genuinely was changing at
that instant. The problem is the delay before you get an answer at all, and that is what the next
section buys time for.

\section{The Double-Flop Synchronizer}
\label{c4:a3}
The standard fix is the \textbf{double-flop synchronizer}:
\begin{itemize}
  \item Place two D flip-flops in series on every asynchronous input: every input that changes with
    no fixed relationship to your clock. A push button, a reset line, a signal arriving from
    another clock domain.
  \item An input that is already clocked by \emph{this} clock is not asynchronous, and must
    \textbf{not} be synchronized. Doing so only delays it, and shifts its timing relative to
    everything around it. \code{seq_detect_mealy} in Chapter~\ref{c8:ch} is exactly that case: its
    \code{din} is synchronous serial data, so it goes straight into the machine while
    \code{reset_n} still gets a synchronizer.
  \item Only ever use the output of the \emph{second} flip-flop elsewhere in the design.
\end{itemize}

\begin{textdrawing}
async_in
   │
   ▼
 [FF1] ──────► [FF2] ──────► synchronized, safe to use
   (may go        (extremely likely
   metastable)     to be stable)
\end{textdrawing}

\noindent The same two flip-flops, in time:

\bookfigure{synchronizer_waveform}{\code{async_in} changes just before a clock edge; \code{s1} may go
metastable for a clock period, and \code{s2} comes out clean one edge later.}{c4:fig:synchronizer}

\noindent\code{FF1} is the one exposed to the asynchronous input, so it is the one that might go
metastable. Its output only has to reach \code{FF2}, and it has an entire clock period to get there.
The probability that a metastable node is still unresolved after a full period is small, and
shrinks the longer it is given, so by the time \code{FF2} samples on the next edge, \code{FF1} has
almost always settled. \code{FF2}'s output is what you build the rest of your synchronous logic on.

This is probabilistic, not absolute, and \secref{c4:a4} says what that means. At the speeds used in
this book, tens of MHz, two stages are the standard, industry-accepted answer.

One rule ties it together: \textbf{every asynchronous input gets its own double-flop chain before
it is used anywhere else.} External reset buttons, push buttons, and later any
\textbf{single-bit} signal crossing from another clock domain.

That last word is load-bearing, and it is the one limit worth learning before you over-apply the
rule. A double-flop chain synchronizes \textbf{one bit}. Put a multi-bit bus through a chain per bit
and each bit is individually safe but they are not safe \emph{together}: each independently
resolves to the old or new value on the same edge, so the receiver can sample a combination that
was never on the bus. A counter stepping from \code{0111} to \code{1000} has every bit changing at
once, and a per-bit synchronizer can hand you \code{1111} or \code{0000}, values the counter never
held.

Buses crossing clock domains need a handshake, a FIFO, or Gray coding, which arranges that only one
bit ever changes at a time. This book never crosses clock domains with a bus, so you will not need
them here. Knowing that the two-flop answer stops at one bit is enough for now.

\section{Timing Intuition: Why ``Almost Certainly'' Is Good Enough}
\label{c4:a4}
A synchronizer does not buy ``no metastability''. It buys metastability that gets a full clock
period to resolve before it can affect anything else, and that is a statistical argument rather
than a proof. The correct posture is ``reasonably certain, backed by an extremely low,
well-characterized failure rate'', not ``provably safe under all conditions'', which is why extra
stages exist for extreme cases instead of two being declared sufficient by law of physics.

The relationship has a standard shape, worth seeing once even though you will not compute with it
here:

\begin{console}
MTBF = e^(t_r / tau) / (T_w * f_clock * f_data)
\end{console}

\noindent Read it for what it says rather than for numbers. \code{t_r} is the resolution time a
stage gets, roughly one clock period; \code{tau} and \code{T_w} are both properties of the
flip-flop, \code{tau} being how fast a metastable value decays and \code{T_w} the width of the
window around the clock edge in which the input can provoke one at all. Adding a stage buys another
\code{t_r} in the exponent, so \textbf{the numerator grows exponentially with every stage you add},
against one clock cycle of added latency. The denominator grows linearly with the clock rate and,
just as importantly, with \code{f_data}: how often the asynchronous input actually changes. A
button pressed once a second and a data line toggling at 10\,MHz differ by seven orders of magnitude
in that term alone, which is why the case for a third stage is a question about the input, not
about the clock.

\section{Synchronizing a Reset Signal: Assert Async, Release Sync}
\label{c4:a5}
Resets deserve a specific pattern, which \code{reset_sync} demonstrates. There is no
\file{reset_sync.vhd} in the repository to open: it is Exercise~\ref{c4:ex:8}, and you write it.
The listing below is the whole of its logic; the \code{library}/\code{use} clauses, the
\code{entity}, the \code{architecture} frame and the declaration of the internal \code{reset_s1_n}
are yours to add around it.

\begin{vhdlcode}
process(clock, reset_n) is
begin
    if (reset_n = '0') then
        reset_s1_n <= '0';
        reset_s2_n <= '0';
    elsif (rising_edge(clock)) then
        reset_s1_n <= '1';
        reset_s2_n <= reset_s1_n;
    end if;
end process;
\end{vhdlcode}

\begin{itemize}
  \item \textbf{Assert asynchronously.} When \code{reset_n} drops, both stages are forced to
    \code{0} immediately, with no dependency on the clock. A reset that waited for a clock edge
    would defeat its own purpose: it must take effect the instant it is needed, clock or no clock.
  \item \textbf{Release synchronously.} When \code{reset_n} returns to \code{1} the reset condition
    does not vanish everywhere at once. It shifts out through the two stages like any other
    asynchronous input, so \code{reset_s2_n}, the signal the design actually uses, releases cleanly
    on a clock edge together with everything else. Without this, different flip-flops could leave
    reset on different cycles and momentarily disagree about system state, the same class of bug
    metastability causes.
\end{itemize}
This is why \code{reset_n} is synchronized by its own chain before anything else depends on it,
including the button synchronizer, whose reset input is \code{reset_s2_n} rather than the raw
\code{reset_n}.

\section{Reusing the Chain for Edge Detection --- and, Incidentally, Debouncing}
\label{c4:a6}
The button needs synchronizing, and from Chapter~\ref{c3:ch} you also need to detect \emph{when} it
is pressed rather than its level. \code{button_sync}, which is Exercise~\ref{c4:ex:9} and so is also
yours to write rather than a file in the repository, gets both from a single chain of
\textbf{three} flip-flops:

\begin{textdrawing}
button_n
   │
   ▼
[FF1] → [FF2] → [FF3]
        │        │
        │        └─ previous value
        └─ current value
\end{textdrawing}

\noindent Flip-flops 1 and 2 are \secref{c4:a3}'s synchronizer, making \code{button_s2_n} the stable
current value. Flip-flop 3 stores last cycle's \code{button_s2_n} as \code{button_s3_n}, the
previous value. Flip-flops 1 and 2 are worth noticing on their own: Chapter~\ref{c5:ch}'s
Exercise~\ref{c5:ex:4} extracts that double-flop chain into a reusable \code{sync} module, leaving
the third flip-flop here as the edge detector's own. Comparing stage 2 against stage 3 gives edge
detection for free:

\begin{center}
\code{button_edge_s2 = (not button_s2_n) and button_s3_n}
\end{center}

\noindent which is true for exactly one clock cycle: the one where the button reads pressed now and
did not the cycle before.

\textbf{Synchronization, on its own, only solves metastability.} It says nothing about a mechanical
button's tendency to bounce: to make and break contact several times over the first millisecond or
so, each bounce looking like a separate press. What keeps this book's exercises from being swamped
by bounce is the \emph{combination} of the chain with a detection window that is slow relative to
the bounce:
\begin{itemize}
  \item In this chapter's CircuitVerse exercises the clock period is a full second. Bounce settles
    in far less than that, so by the next edge \code{button_s2_n} reads one clean level. The chain
    looks like it is debouncing, but it is really just sampling too infrequently to see bounce.
  \item On the FPGA the clock is \code{50 MHz}, one sample every \code{20 ns}, while bounce lasts
    from under a millisecond to a few tens of milliseconds, many thousands of cycles. At that speed
    \code{button_sync} faithfully detects \emph{every} bounce transition as its own edge, and the LED
    could toggle several times for one physical press.
\end{itemize}
So precisely: the double-flop synchronizer reliably solves metastability at the speeds used here.
Whether it also \emph{looks} like debouncing depends entirely on the clock period being slow
relative to the switch's bounce time, which is true in CircuitVerse and not automatically true on
the FPGA.

\section{What Real Hardware Debouncing Adds}
\label{c4:a7}
This book has no dedicated debounce chapter, and the worked examples deliberately keep using
\secref{c4:a6}'s circuit, bounce and all. It is good enough for a demo LED. A production design adds
one of the following \emph{on top}:
\begin{itemize}
  \item \textbf{An analog RC low-pass filter}, optionally with a Schmitt-trigger buffer, on the
    button input before it reaches an FPGA pin. The RC pair smooths the fast make/break transitions
    into a single gradual edge, and the Schmitt trigger, with its two switching thresholds, turns
    that back into one clean digital transition. This solves bounce before the signal is even
    digital, but costs board components and does not help for inputs that are not mechanical
    switches.
  \item \textbf{A digital debounce counter.} Rather than acting on the next detected edge, require
    the input to hold a new level for some minimum duration, a few milliseconds' worth of clock
    cycles, before accepting it. That builds your own artificially slow detection window on-chip,
    the same effect the 1-second CircuitVerse clock gives for free, without slowing the system
    clock. You will meet counters in Chapter~\ref{c6:ch} and timers, which this technique needs, in
    Chapter~\ref{c7:ch}.
\end{itemize}
Both sit on top of the synchronizer, not instead of it: even a perfectly debounced signal is still
asynchronous to your clock and still needs a double-flop chain.

\section{Generics: One Module, Several Sizes}
\label{c4:a8}
\code{reset_sync} and \code{button_sync} are built from the same idea, but only one has a question
to answer. A reset is one bit and always will be. How many \emph{buttons} does a design have? This
chapter's \code{led_toggle_sync} has one, Chapter~\ref{c8:ch}'s \code{fsm_led} has two. Writing the
same three-flop chain out once per width is exactly the duplication that splitting them into
modules set out to avoid.

A \textbf{generic} answers it: a value an entity declares alongside its ports, fixed when the design
is built rather than changing while it runs.

\begin{vhdlcode}
entity button_sync is
    generic(COUNT: natural range 1 to 3 := 1);
    port(clock, reset_s2_n: in  std_logic;
         button_n         : in  std_logic_vector(COUNT-1 downto 0);
         button_edge_s2   : out std_logic_vector(COUNT-1 downto 0));
end entity;
\end{vhdlcode}

\begin{itemize}
  \item \code{COUNT} behaves like a constant inside the architecture. It cannot change while the
    circuit runs; it is decided before a single gate is placed.
  \item \code{:= 1} is its default, used by any instantiation that says nothing.
  \item The port widths are written \emph{in terms of it}, and so is every internal signal that
    follows.
\end{itemize}

\subsection{Overriding it}
\label{c4:override}
An instantiation passes generics with a \code{generic map}, before the \code{port map}:

\begin{vhdlcode}
button_sync1: entity work.button_sync
    generic map(2)
    port map(clock, reset_s2_n, button_n, button_edge_s2);
\end{vhdlcode}

\noindent Generics are matched \textbf{by position}, in declaration order, exactly as ports are in
\secref{c2:a6}. Leaving the \code{generic map} out applies the default; this book writes it
explicitly even when it matches, because a reader should not have to open the module to find out
how wide an instance is.

\subsection{When a generic is worth having}
\label{c4:genericworth}
\code{button_sync} has one and \code{reset_sync} has none, and that difference is the whole rule:
\textbf{add a generic when something genuinely varies between instances, not on principle.}

\code{button_sync} is instantiated at \code{COUNT = 1} by \secref{c4:a9}'s \code{led_toggle_sync}
and at \code{COUNT = 2} by Chapter~\ref{c8:ch}'s \code{fsm_led}: two widths, one file.
\code{reset_sync} is one bit everywhere, so a \code{WIDTH} generic on it would be a parameter nobody
ever passes anything but \code{1} to, and one more thing for a reader to check. That second
instantiation is what makes a generic worth having; a parameter with exactly one value in the
entire book would just be indirection. \code{button_sync}'s own testbench instantiates it at
\emph{both} widths in one simulation, precisely to prove that one file really does serve both.

You will meet one more generic in Chapter~\ref{c7:ch}: \code{timer}'s \code{TICK_COUNT}, which is a
count rather than a width, and is what lets the same file measure a fiftieth of a second in one
design and a whole second in another.

\section{Worked Example: \code{led_toggle_sync}}
\label{c4:a9}
This ties \secref{c4:a1} to \secref{c4:a6} together into Chapter~\ref{c3:ch}'s toggle circuit, made
safe for a real, asynchronous button and reset.

\begin{booktable}{@{}p{6em}p{5em}L@{}}
  \thead{Port} & \thead{Direction} & \thead{Description} \\
  \midrule
  \code{clock} & in & \code{50 MHz} system clock on the DE0-CV board. \\
  \code{reset_n} & in & Active-low asynchronous reset, wired to a push button. \\
  \code{button_n} & in & Active-low push button; toggles the LED on its falling edge. \\
  \code{led} & out & LED, toggled on each synchronized, edge-detected button press. \\
\end{booktable}

\noindent Internal signals: \code{led_s} holds the LED state, \code{reset_s2_n} is the reset after
the two-flop synchronizer (the \code{s2} suffix is this chapter's convention for ``synchronized
with two flip-flops''), and \code{button_edge_s2} is the one-cycle pulse marking a synchronized
falling edge on \code{button_n}.

Like Chapter~\ref{c2:ch}'s two-digit display, this design is built from more than one entity, using
the instantiation from \secref{c2:a6} and the generic from \secref{c4:a8}. The difference is what
the pieces do: the display instantiated one module twice to do the same job on two halves of its
input, while this instantiates two \emph{different} modules with one job each. Both are the same
mechanism. Composing a design out of modules is not a special technique for repeated logic; it is
how designs are built once they outgrow a single file.

The two chains live in two subcomponents:
\begin{itemize}
  \item \code{reset_sync} is \secref{c4:a5}'s two-flop reset synchronizer. It takes the raw
    \code{reset_n} and returns \code{reset_s2_n}.
  \item \code{button_sync} is \secref{c4:a6}'s three-flop synchronizer-plus-edge-detector. It takes
    the already synchronized \code{reset_s2_n} and pulses \code{button_edge_s2} once per press.
\end{itemize}
Keeping them apart matters more than it looks. A design with no button of its own still needs its
reset synchronized, and with two modules it simply instantiates the one it wants, as
Chapter~\ref{c8:ch}'s \code{seq_detect_mealy} does. Welded into one module, that design would have
to instantiate the whole thing and disable the half it does not use.

\code{button_sync} takes a \emph{vector} of buttons so one file can serve a two-button design. This
design has one, so \file{led_toggle_sync.vhd} declares a pair of one-element vectors to connect to
it:

\begin{vhdlcode}
signal button_n_v, button_edge_s2_v: std_logic_vector(0 downto 0);
...
button_n_v(0)  <= button_n;
button_edge_s2 <= button_edge_s2_v(0);
\end{vhdlcode}

\noindent That is plumbing, not logic: two wires renamed.

The top-level architecture contains nothing but the toggle process itself, and it only ever touches
the already-synchronized \code{reset_s2_n} and \code{button_edge_s2}, never the raw inputs:

\begin{vhdlcode}
LED_PROCESS: process(clock, reset_s2_n) is
begin
    if (reset_s2_n = '0') then
        led_s <= '0';
    elsif (rising_edge(clock)) then
        if (button_edge_s2 = '1') then
            led_s <= not led_s;
        end if;
    end if;
end process;
\end{vhdlcode}

\noindent The full source of the top level is \file{lectures/L04/led_toggle_sync/led_toggle_sync.vhd}:

\vhdlfile{lectures/L04/led_toggle_sync/led_toggle_sync.vhd}

\noindent The two subcomponents are not in the repository, because writing them is
Exercises~\ref{c4:ex:8} and~\ref{c4:ex:9}. What this chapter gives you of each is:
\code{reset_sync}'s process in \secref{c4:a5} above, and \code{button_sync}'s chain as a diagram and
an equation in \secref{c4:a6} with its entity in \secref{c4:a8}. The surrounding entity and
architecture in the first case, and the whole architecture body in the second, are the exercises'
own work. Copy your own into \file{led_toggle_sync/} to build it.

\textbf{By hand, in CircuitVerse first.} Before touching VHDL, build it as a gate network: two D
flip-flops in series for the reset synchronizer; two for \code{button_n} plus a third storing the
previous synchronized value; an AND gate combining the inverted current value with the previous one
to produce \code{button_edge_s2}; and a toggle flip-flop for \code{led_s}, its inverted output fed
back to its input, enabled only when \code{button_edge_s2 = 1}. Set the clock period to
\mbox{\code{1000 ms}} so you can watch each stage by eye.

\bookimage{1.0}{lectures/L04/appendix/images/led_toggle_sync_circuit.png}{The synchronized LED
toggle in CircuitVerse: the reset and button double-flop synchronizers, the edge-detecting AND gate,
and a toggle flip-flop for \code{led_s}.}{c4:fig:circuit}

\noindent One thing in the drawing needs reading carefully: the flip-flops carry a \code{preset}
input, and on the three button-synchronizer stages it is tied to \code{1}. That is not decoration.
In CircuitVerse, \code{preset} is the value a flip-flop loads when its reset is asserted, so tying
it to \code{1} is what makes \code{button_s1_n}, \code{button_s2_n} and \code{button_s3_n} come out
of reset holding \code{1} rather than \code{0}.

That matches \code{button_sync}, which resets the same three stages to \code{(others => '1')}, and
it matters for the same reason: the buttons are active-low, so \code{1} means ``released''. Reset
the chain to \code{0} instead and it comes up claiming the button is already held down. The reset
flip-flops leave \code{preset} alone, because there a reset to \code{0} really is all the circuit
asks for.

Press the button in the simulation and watch it propagate. The LED does not toggle instantly: it
waits for the synchronized, edge-detected pulse to reach the toggle flip-flop, which takes a small,
fixed number of clock pulses, exactly the latency \secref{c4:a3} called the price of safety.

\textbf{On the FPGA:} the course demonstrates the same design on the Terasic DE0-CV (device
\code{5CEBA4F23C7N}), with \code{clock} on the \code{50 MHz} oscillator, \code{reset_n} and
\code{button_n} on two push buttons, and \code{led} on an LED; Appendix~\ref{app:quartus} is the
workflow, for a reader who has a board. The LED toggles on every press, and per \secref{c4:a6} a
real press may occasionally produce more than one toggle from contact bounce.
